\section{A shorter finite UW expression}
\label{sec:comp-UW-simplified}
The finite $UW$ contribution already has the integrated endpoint form
\eqref{eq:comp-real-UW-final-primitive}. Its thirteen inner monomials
contain repeated weight-two functions. Combining them before the
rational primitive reduces the endpoint expression. Throughout this
section $H=s-1>0$, the electron mass is one, and $0<t<1$.
Define
\begin{equation}
 P=1+Ht-Ht^2,\qquad \ell=\ln(1+Ht),\qquad
 p=\ln P,\qquad b=\ln(1+Ht/2).
\end{equation}
All three logarithms are real. The two weight-two combinations needed
below are
\begin{align}
 J_2(H,t)&=\operatorname{Li}_2(-1-Ht)+\frac{\pi^2}{12}
       +\ell\ln(2+Ht),\\
 J(H,t)&=\operatorname{Li}_2(t/P)-\operatorname{Li}_2(t)
       -p\ln(1-t)+\frac{p^2}{2}.
 \label{eq:comp-UW-J-short}
\end{align}
The dilogarithm arguments in $J$ obey $0<t/P\leq t<1$.
The argument in $J_2$ is negative. These expressions therefore fix the
physical real branches directly.

\paragraph{Exact coefficient combinations.} Let $r_f(H,t)$ denote the
coefficient of the atom monomial $f$ in the thirteen-term inner result.
Its rational coefficients obey
\begin{equation}
 r_{d_3}=-r_{\ell\ln t}=r_{\ell p},\qquad
 r_{d_2}=r_{\ell\ln(2+Ht)},\qquad
 [\ln2]r_1=-r_{\ln(2+Ht)}.
 \label{eq:comp-UW-coefficient-combinations}
\end{equation}
The last notation means the coefficient of $\ln2$ in $r_1$.
These are polynomial identities after multiplication by their common
rational denominators, for arbitrary $H$ and $t$. Here
$d_2=\operatorname{Li}_2(-1-Ht)+\pi^2/12$ and $d_3$ is the two-dilogarithm
difference defined in the inner integration. Thus its finite kernel is
\begin{equation}
 \Phi_{UW}=A_0+A_1\ell+A_2\ell^2+A_p p+A_b b
             +A_d\operatorname{Li}_2(-Ht)+A_{J_2}J_2+A_JJ,
 \label{eq:comp-UW-eight-kernel}
\end{equation}
with the following factored sums of the original rational coefficients:
\begin{align}
 A_0&=\left.r_1\right|_{\ln2=0}+r_{\ln H}\ln H,
 &A_1&=r_\ell+r_{\ell\ln H}\ln H,\\
 A_2&=r_{\ell^2},& A_p&=r_p,\\
 A_b&=r_{\ln(2+Ht)},& A_d&=r_{d_1},\\
 A_{J_2}&=r_{d_2},& A_J&=r_{d_3}.
\end{align}
Here $d_1=\operatorname{Li}_2(-Ht)$. Keeping these sums factored avoids
expanding large common numerators. The explicit eight coefficients are
stored with the calculation in
\texttt{\detokenize{uw_simplified_kernel.json}}.

\paragraph{Dilogarithm reduction.} To obtain $J$ without guessing a
branch, first apply the negative-real-axis identity
$\operatorname{Li}_2(-z)+\operatorname{Li}_2(-1/z)
=-\pi^2/6-\ln^2z/2$, with $z>0$, to $d_3$.
It gives the analytic combination $D=d_3-\ell\ln t$ derived at the lower
endpoint. Next use
\begin{equation}
 \operatorname{Li}_2(-z)
 =-\operatorname{Li}_2\!\left(\frac{z}{1+z}\right)
  -\frac12\ln^2(1+z),\qquad z>0.
\end{equation}
The two positive arguments are
$t/[(1-t)(1+Ht)]$ and $t/(1-t)$.
Their transformed arguments are $t/P$ and $t$. The logarithms combine
exactly to give $J=D+\ell p$, hence \eqref{eq:comp-UW-J-short}.
In particular no explicit $\ln t$ survives. Both $J(0)=0$ and $J_2(0)=0$.
At the upper endpoint their analyticity follows from the original
$d_3-\ell\ln t+\ell p$ and $d_2+\ell\ln(2+Ht)$ representations.
This observation also fixes the value at $t=1$, where separate terms
of \eqref{eq:comp-UW-J-short} have logarithmic derivatives.

\paragraph{The integrated result.} Introduce the letters
\begin{equation}
 a=-\frac1H,\qquad c=-\frac2H,\qquad
 r_\pm=\frac{1\pm\sqrt{1+4/H}}2.
\end{equation}
In the same GPL convention as the rational primitive,
\begin{align}
 J_2(H,t)&=G(c,a;t),\\
 J(H,t)&=\sum_{r=r_\pm}
       \bigl[G(0,r;t)-G(1,r;t)+G(r,a;t)\bigr]-G(0,a;t).
\end{align}
For example $J_2'=H\ell/(2+Ht)$ and $J_2(0)=0$ prove its single-word
form. The formula for $J$ follows from the shuffle product
$G(a;t)G(r;t)=G(a,r;t)+G(r,a;t)$ and the expression for $D$.
The complete finite integral is therefore
\begin{equation}
\begin{aligned}
 \mathcal F_{UW}(H)=\lim_{Z\to1^-}\bigg\{&
 \mathfrak T[A_0;\varnothing](Z)+\mathfrak T[A_1;a](Z)
 +2\mathfrak T[A_2;a,a](Z)\\
 &+\mathfrak T[A_b;c](Z)+\mathfrak T[A_{J_2};c,a](Z)
 -\mathfrak T[A_d+A_J;0,a](Z)\\
 &+\sum_{r=r_\pm}\Bigl(
 \mathfrak T[A_p;r](Z)+\mathfrak T[A_J;0,r](Z)\\
 &\hspace{26mm}-\mathfrak T[A_J;1,r](Z)
 +\mathfrak T[A_J;r,a](Z)\Bigr)\bigg\}.
\end{aligned}
\label{eq:comp-UW-short-primitive}
\end{equation}
The operator $\mathfrak T$ is the finite rational-to-GPL primitive
already evaluated by the terminating rules above, including its
zero finite part at the lower endpoint. Thus this is a shorter
representation of the completed integral, with fourteen primitive
contributions. It is not a replacement by an unevaluated quadrature.
The terms with first letter one must be summed before taking the
upper limit. The total cross section uses this sector with the same
factor $1/(8sH)$ and the second Bose-related contribution as before.

\paragraph{Removal of apparent endpoint poles.} The certificates in
Sections~\ref{sec:comp-UW-upper-certificate} and
\ref{sec:comp-UW-lower-certificate} also give a form in which each
summand is finite at both endpoints. For any one of the eight functions
$f_i$ in \eqref{eq:comp-UW-eight-kernel}, decompose its coefficient as
\begin{equation}
 A_i(t)=A_i^{\mathrm{reg}}(t)
       +\sum_{e=0,1}\sum_{k=1}^{m_{ie}}\frac{c_{iek}}{(t-e)^k}.
\end{equation}
The regular coefficient has no pole at zero or one. Since every $f_i$
is analytic there, define its finite Taylor subtraction
\begin{equation}
 f_i^{[e,k]}(t)=\frac{f_i(t)-\displaystyle\sum_{n=0}^{k-1}
                  f_i^{(n)}(e)(t-e)^n/n!}{(t-e)^k}.
\end{equation}
Then the same finite kernel satisfies the exact identity
\begin{equation}
 \Phi_{UW}(t)=\sum_i A_i^{\mathrm{reg}}(t)f_i(t)
          +\sum_{i,e,k}c_{iek}f_i^{[e,k]}(t).
\end{equation}
Indeed, the omitted Taylor pieces are exactly the principal parts
of the full kernel at zero and one, which the two symbolic certificates
show to vanish. These divided differences are explicit logarithms and
dilogarithms minus finite polynomials. They introduce no new integral.
The factored form \eqref{eq:comp-UW-short-primitive} is shorter for
presenting the integrated result.
